\documentclass[11pt,a4paper]{article}

% ------------------------------------------------------------------
%  RUNAWAY : Operation Eclipse - Game Design Document v4.0
%  Genere automatiquement a partir du code source reellement implemente.
% ------------------------------------------------------------------
\usepackage[utf8]{inputenc}
\usepackage[T1]{fontenc}
\usepackage[french]{babel}
\usepackage[a4paper,margin=2.3cm]{geometry}
\usepackage{amsmath,amssymb}
\usepackage{xcolor}
\usepackage{booktabs}
\usepackage{array}
\usepackage{longtable}
\usepackage{enumitem}
\usepackage{fancyhdr}
\usepackage{titlesec}
\usepackage[hidelinks,colorlinks=true,linkcolor=eclipsered,urlcolor=eclipseblue]{hyperref}

\definecolor{eclipsered}{RGB}{192,57,43}
\definecolor{eclipseblue}{RGB}{90,120,200}
\definecolor{eclipsedark}{RGB}{30,30,40}
\definecolor{codebg}{RGB}{244,244,248}

\titleformat{\section}{\Large\bfseries\color{eclipsered}}{\thesection.}{0.5em}{}
\titleformat{\subsection}{\large\bfseries\color{eclipsedark}}{\thesubsection}{0.5em}{}

\setlength{\headheight}{14pt}
\pagestyle{fancy}
\fancyhf{}
\lhead{\textcolor{eclipsered}{\textbf{RUNAWAY}} -- Operation Eclipse}
\rhead{GDD v4.0}
\cfoot{\thepage}
\renewcommand{\headrulewidth}{0.4pt}

\newcommand{\code}[1]{\colorbox{codebg}{\texttt{#1}}}
\newcommand{\param}[2]{\code{#1} $=$ #2}

\begin{document}

% ============================== PAGE DE TITRE ==============================
\begin{titlepage}
\centering
\vspace*{3cm}
{\Huge\bfseries\color{eclipsered} RUNAWAY\par}
\vspace{0.5cm}
{\LARGE Operation Eclipse\par}
\vspace{2cm}
{\Large\bfseries GAME DESIGN DOCUMENT\par}
\vspace{0.3cm}
{\large VERSION 4.0 --- Edition production / mobile\par}
\vspace{2cm}
\begin{minipage}{0.8\textwidth}\centering\itshape
« Tu entres. Tu braques. Tu cours. L'Organisation Fant\^ome ne te laissera
pas partir facilement. »
\end{minipage}
\vfill
{\small Jeu d'action / infiltration top-down --- Godot 4.6.2 / GDScript\par}
{\small Campagne monde ouvert continu en 3 environnements --- Cible : Android (Play Store) \& Desktop\par}
\vspace{0.5cm}
{\small Document mis \`a jour le \today\par}
\end{titlepage}

\tableofcontents
\newpage

% ============================== 1. VISION ==============================
\section{Vision du Jeu}

\subsection{1.1 Concept}
RUNAWAY est un jeu d'action top-down dans lequel le joueur incarne \textbf{Agent
Z\'ero}, un espion infiltr\'e qui doit traverser une s\'erie de zones ennemies,
r\'ecup\'erer des donn\'ees classifi\'ees et s'exfiltrer avant que l'alarme ne se
g\'en\'eralise. La cam\'era suit le joueur en vue de dessus. Le jeu m\^ele action
directe (tir au pistolet silencieux) et furtivit\'e (couverture automatique,
\'evitement des c\^ones de vision).

\textbf{\'Evolution majeure depuis la v3 :} le jeu n'est plus un niveau unique
rejouable mais une \textbf{campagne « monde ouvert continu » en trois
environnements encha\^in\'es}. La sortie d'une zone devient l'entr\'ee de la
suivante~: l'espion sort du centre de recherche, se retrouve dans une for\^et,
puis dans une ville. Le score, les vies et le temps sont \textbf{cumul\'es} d'un
bout \`a l'autre de la campagne.

\subsection{1.2 Boucle de Jeu Core}
\begin{enumerate}[leftmargin=*]
  \item \textbf{Infiltrer} --- entrer dans la zone en \'evitant les c\^ones de vision.
  \item \textbf{Agir} --- accomplir l'objectif local (pirater le coffre au niveau~0)
        ou traverser un point de passage (pont sur la rivi\`ere / le fleuve).
  \item \textbf{Fuir} --- rejoindre le point d'exfiltration sous la pression des gardes.
  \item \textbf{Encha\^iner} --- la sortie recharge la sc\`ene sur l'environnement suivant,
        avec bonus d'\'etape, jusqu'\`a l'exfiltration finale.
\end{enumerate}

\subsection{1.3 Piliers de Design}
\begin{itemize}[leftmargin=*]
  \item \textbf{Tension permanente} --- l'alarme monte et redescend ; rester discret est r\'ecompens\'e.
  \item \textbf{IA cr\'edible} --- les gardes voient, entendent, communiquent, encerclent et anticipent.
  \item \textbf{Lisibilit\'e} --- c\^ones de vision color\'es, HUD minimal, minimap.
  \item \textbf{Accessibilit\'e multi-support} --- clavier/souris \emph{et} double-stick tactile.
  \item \textbf{P\'erim\`etre ma\^itris\'e} --- production compl\`ete et stable plut\^ot qu'ambition diffuse.
\end{itemize}

% ============================== 2. UNIVERS ==============================
\section{Univers et Narration}

\subsection{2.1 Contexte}
2031. L'Organisation Fant\^ome contr\^ole des gouvernements via des donn\'ees
compromettantes stock\'ees physiquement. Agent Z\'ero, dernier op\'erationnel de
l'unit\'e ECLIPSE, doit r\'ecup\'erer ces preuves et fuir \`a travers le dispositif
de s\'ecurit\'e ennemi : un complexe de recherche, la for\^et qui l'entoure, puis
la ville o\`u l'attend l'exfiltration finale.

\subsection{2.2 Agent Z\'ero --- Fiche Personnage}
Espion solitaire, \'equip\'e d'un pistolet silencieux (chargeur de 8 balles),
capable de sprinter sur de courtes dur\'ees et de se mettre \`a couvert
automatiquement. Trois points de vie, invincibilit\'e br\`eve apr\`es un coup.

\subsection{2.3 L'Organisation Fant\^ome}
\begin{itemize}[leftmargin=*]
  \item R\'eseau criminel transnational, infiltr\'e dans les gouvernements.
  \item Dispose de gardes arm\'es de trois profils et d'une coordination radio.
  \item Point faible : les donn\'ees sont stock\'ees physiquement (pas de piratage \`a distance).
\end{itemize}

\subsection{2.4 Ton Narratif}
Pas de cin\'ematiques. La narration passe par l'ambiance : musique adaptative
selon l'\'etat d'alarme, environnements sombres, gardes qui se pr\'eviennent
mutuellement. La pression se \emph{ressent}.

\subsection{2.5 Structure de Campagne}
La campagne est un encha\^inement continu de trois environnements
(constante \param{LEVEL\_COUNT}{3}). La progression est port\'ee par le singleton
\code{GameState} qui survit au rechargement de sc\`ene.

\begin{center}
\begin{tabular}{clll}
\toprule
\textbf{Niv.} & \textbf{Environnement} & \textbf{Objectif local} & \textbf{Sortie} \\
\midrule
0 & Centre de recherche & Pirater le coffre (donn\'ees) & Exfil haut-gauche \\
1 & La For\^et & --- (traversée) & Franchir le pont, rejoindre l'exfil \\
2 & La Ville & --- (traversée) & Pont vers l'h\'elipad d'extraction \\
\bottomrule
\end{tabular}
\end{center}

Seul le niveau~0 exige un objectif (\code{needs\_objective = true}) ; les niveaux~1
et~2 sont des travers
\'ees o\`u atteindre la sortie suffit. Chaque \'etape franchie octroie un bonus de
$+300$~pts.

% ============================== 3. MECANIQUES ==============================
\section{M\'ecaniques de Jeu}

\subsection{3.1 Contr\^oles}
Deux sch\'emas coexistent. Le sch\'ema tactile s'active \emph{paresseusement} :
l'overlay reste invisible et ne perturbe pas le clavier tant qu'aucun doigt n'a
touch\'e l'\'ecran.

\begin{longtable}{p{3.2cm} p{4.3cm} p{6.5cm}}
\toprule
\textbf{Action} & \textbf{Clavier/Souris} & \textbf{Tactile (Android)} \\
\midrule
\endhead
D\'eplacement & WASD / Fl\`eches & Stick gauche (moiti\'e gauche de l'\'ecran) \\
Sprint & Maj & Stick gauche pouss\'e \`a fond ($\geq 82\,\%$) \\
Viser & Orientation auto vers le d\'eplacement & Stick droit (direction) \\
Tirer & Clic gauche & Stick droit maintenu \\
Recharger & R & Bouton \textbf{R} (bord droit) \\
Interagir & E & Bouton \textbf{E} (bord droit) \\
Pause & \'Echap / P & Bouton \textbf{Retour} Android \\
\bottomrule
\end{longtable}

\subsection{3.2 Syst\`eme de D\'eplacement}
Vitesse de base \param{SPEED}{270} px/s. Le sprint multiplie la vitesse
(\param{SPRINT\_MULT}{1.6}) et consomme une r\'eserve d'\'energie :

\[
\dot{e}(t) =
\begin{cases}
-1 & \text{si sprint actif } (e > 0,\ \vec{v}\neq 0,\ \text{hors couverture})\\[2pt]
+\,\texttt{SPRINT\_REGEN} & \text{sinon}
\end{cases}
\qquad e \in [0,\ \texttt{SPRINT\_MAX}]
\]

avec \param{SPRINT\_MAX}{2.2}~s et \param{SPRINT\_REGEN}{0.8}~s$^{-1}$.
En couverture, la vitesse est r\'eduite : $v = \texttt{SPEED}\cdot\texttt{COVER\_SPEED\_MULT}$,
\param{COVER\_SPEED\_MULT}{0.7}.

\subsection{3.3 Syst\`eme de Tir}
\begin{itemize}[leftmargin=*]
  \item Chargeur : \param{MAG\_SIZE}{8} balles ; rechargement \param{RELOAD\_TIME}{1.2}~s (vuln\'erable).
  \item Cadence : \param{FIRE\_COOLDOWN}{0.22}~s entre deux tirs.
  \item Projectile : vitesse $700$~px/s, dur\'ee de vie $1.5$~s, 1 d\'eg\^at.
  \item Tir silencieux : n'alerte pas les gardes hors champ de vision.
\end{itemize}

\subsection{3.4 Couverture Automatique}
D\`es qu'un mur/obstacle est d\'etect\'e par la zone \code{CoverArea}, le joueur
passe en couverture : hitbox r\'eduite \`a \param{COVER\_HITBOX\_MULT}{0.6} (-40\,\%)
et vitesse \`a $0.7$ (-30\,\%).

\subsection{3.5 Syst\`eme d'Alarme}
L'alarme est une machine \`a trois \'etats globaux pilotant la musique et le score.

\begin{center}
\begin{tabular}{lll}
\toprule
\textbf{\'Etat} & \textbf{Mont\'ee} & \textbf{Redescente} \\
\midrule
CALME & \'etat initial & --- \\
SUSPICION & un garde voit/entend bri\`evement & apr\`es \texttt{SUSPICION\_TIMEOUT}$=5$\,s sans contact \\
ALERTE & combat confirm\'e ou tir entendu & apr\`es \texttt{ALERT\_TIMEOUT}$=7$\,s sans combat $\to$ SUSPICION \\
\bottomrule
\end{tabular}
\end{center}

Formellement, deux minuteries d\'ecroissent dans \code{GameState.\_process} :
\[
t_{\text{susp}} \mathrel{-}= \Delta t \;\Rightarrow\; \text{CALME quand } t_{\text{susp}}\leq 0,
\qquad
t_{\text{combat}} \mathrel{+}= \Delta t \;\Rightarrow\; \text{SUSPICION quand } t_{\text{combat}}\geq 7.
\]
Chaque frame de combat d'un garde appelle \code{notify\_combat()} qui remet
$t_{\text{combat}}\!=\!0$, maintenant l'alerte tant que le joueur est engag\'e.

\subsection{3.6 Syst\`eme de Score}
Le score est \textbf{cumulatif sur toute la campagne}. Bar\`eme :

\begin{center}
\begin{tabular}{lr}
\toprule
\textbf{\'Ev\'enement} & \textbf{Points} \\
\midrule
Objectif principal (coffre) & $+1000$ \\
Garde neutralis\'e & $+100$ /garde \\
\'Etape de campagne franchie & $+300$ /\'etape \\
Exfiltration finale r\'eussie & $+500$ \\
Discr\'etion totale (alarme jamais d\'eclench\'ee) & $+500$ \\
Point de vie perdu & $-50$ /PV \\
Temps de mission & $-1$ /seconde \\
\bottomrule
\end{tabular}
\end{center}

Le score final retranche le temps une seule fois :
\[
\texttt{final\_score} = \texttt{score} - \big\lfloor \texttt{mission\_time} \big\rfloor .
\]

\subsection{3.7 Sant\'e et Invincibilit\'e}
\param{MAX\_HEALTH}{3}. Apr\`es un coup, le joueur est invincible
\param{IFRAME\_TIME}{0.5}~s (clignotement visuel). \`A $0$~PV : \'echec de mission.

% ============================== 4. IA ==============================
\section{Intelligence Artificielle des Ennemis}
C'est le c\oe ur tactique du jeu. Cette section d\'ecrit l'IA \emph{r\'eellement
impl\'ement\'ee} dans \code{guard.gd}, formules \`a l'appui.

\subsection{4.1 Types de Gardes}
\begin{center}
\begin{tabular}{llcc}
\toprule
\textbf{Type} & \textbf{Comportement} & \textbf{Arm\'e} & \textbf{PV} \\
\midrule
BASIC & Patrouille, recule et appelle des renforts & Non & 2 \\
ALERT & Cherche activement, tire & Oui & 2 \\
ARMED & Tire, encercle, tient la distance id\'eale & Oui & 3 \\
\bottomrule
\end{tabular}
\end{center}

\subsection{4.2 Machine \`a \'Etats}
Quatre \'etats : \textbf{PATROL} $\to$ \textbf{SUSPICION} $\to$ \textbf{COMBAT}
$\to$ \textbf{SEARCH}, avec retours.

\begin{center}\small
\begin{tabular}{lll}
\toprule
\textbf{De} & \textbf{Vers} & \textbf{Condition} \\
\midrule
PATROL & SUSPICION & voit ou entend le joueur \\
SUSPICION & COMBAT & contact visuel confirm\'e ($t_{\text{vu}}\geq$ \texttt{CONFIRM\_TIME}) \\
SUSPICION & PATROL & a fouill\'e le dernier point vu $2.5$\,s \\
COMBAT & SEARCH & perd le contact visuel \\
SEARCH & COMBAT & retrouve le joueur (vue ou ou\"ie) \\
SEARCH & PATROL & fouille infructueuse pendant \texttt{SEARCH\_TIME}$=6$\,s \\
\bottomrule
\end{tabular}
\end{center}

\subsection{4.3 Perception : Champ de Vision}
C\^one de demi-angle \param{VISION\_ANGLE}{$45^\circ$} (90$^\circ$ total). Port\'ee
\param{RANGE\_CALM}{180} px en patrouille, \param{RANGE\_ALERT}{260} px en alerte.
Un \code{RayCast2D} sur la couche \emph{world} bloque la vision derri\`ere les
murs, caisses et l'eau.

\textbf{Vision p\'eriph\'erique (nouveaut\'e).} La d\'etection n'est plus uniforme :
nette dans l'axe, sa port\'ee d\'ecro\^it vers les bords du c\^one. Pour un \'ecart
angulaire $\theta$ \`a la ligne de vis\'ee et $\theta_{1/2}=45^\circ$ :
\[
R_{\text{eff}}(\theta) = R \cdot \Big(1 - 0.45\,\tfrac{\theta}{\theta_{1/2}}\Big),
\qquad 0 \leq \theta \leq \theta_{1/2}.
\]
Le joueur est vu si $d \leq R_{\text{eff}}(\theta)$ et la ligne est d\'egag\'ee.
Au centre la port\'ee est de $100\,\%$, au bord du c\^one elle tombe \`a $55\,\%$.

\textbf{Sens de proximit\'e.} Un intrus coll\'e ($d \leq$ \texttt{PROXIMITY\_SENSE}$=48$\,px)
est rep\'er\'e m\^eme hors du c\^one.

\textbf{Ou\"ie.} Le sprint du joueur g\'en\`ere du bruit ; un garde l'entend si
$d \leq$ \texttt{NOISE\_HEARING}$=120$\,px.

\subsection{4.4 Confirmation de D\'etection}
La confirmation est plus rapide de pr\`es, plus lente de loin. Le compteur $t_{\text{vu}}$
s'accumule avec un facteur de distance :
\[
\dot{t}_{\text{vu}} = \Delta t \cdot \operatorname{clamp}\!\Big(\tfrac{\texttt{RANGE\_ALERT}}{\max(d,1)},\ 0.55,\ 2.2\Big),
\qquad \text{COMBAT d\`es } t_{\text{vu}} \geq \texttt{CONFIRM\_TIME}=0.55\,\text{s}.
\]

\subsection{4.5 Comportement de Combat : Encerclement}
Les gardes arm\'es ne s'agglutinent plus face au joueur : chacun occupe un
\textbf{secteur distinct} autour de lui. L'angle d'encerclement est stable et
unique par garde, tir\'e d'une \emph{suite de Weyl} sur l'identifiant d'instance :
\[
\varphi_i = 2\pi \cdot \operatorname{frac}\!\big(\,\texttt{id}_i \cdot 0.6180339887\,\big)
\quad(\text{nombre d'or}),
\]
ce qui r\'epartit les bearings de mani\`ere quasi-uniforme. Le garde vise le poste
\[
\vec{p}_{\text{poste}} = \vec{p}_{\text{joueur}} + R_{\text{pref}}\,(\cos\varphi',\ \sin\varphi'),
\quad \varphi' = \varphi_i + 0.22\,s,\ s\in\{-1,+1\}\ \text{(strafe)},
\]
avec la distance de tir id\'eale \param{PREFERRED\_RANGE}{190}\,px et une zone morte
\param{RANGE\_HYSTERESIS}{45}\,px pour \'eviter le yo-yo. Il rejoint son poste via
le pathfinding (contournement des murs) tout en gardant l'arme braqu\'ee sur la cible.

Le garde \textbf{non arm\'e} (BASIC) se replie en biais \emph{vers l'alli\'e le plus
proche} (regroupement) et appelle des renforts toutes les $1.2$\,s.

\subsection{4.6 Vis\'ee Anticip\'ee (lead targeting)}
Le garde vise l\`a o\`u le joueur \emph{sera}, compte tenu du temps de vol de la
balle ($v_b = 700$\,px/s). On r\'esout par point fixe (2 it\'erations) :
\[
\vec{a}_0 = \vec{p}_{\text{joueur}},\qquad
\vec{a}_{k+1} = \vec{p}_{\text{joueur}} + \vec{v}_{\text{joueur}}\cdot
\frac{\lVert \vec{a}_k - \vec{p}_{\text{garde}} \rVert}{v_b}.
\]
$\vec{a}_2$ est le point vis\'e. Fuir en ligne droite devient dangereux.

\subsection{4.7 \'Evitement du Tir Ami}
Avant de tirer, le garde v\'erifie qu'aucun alli\'e n'est dans son couloir de tir.
Soit $\vec{d}$ la direction normalis\'ee vers le point vis\'e et $D$ la distance \`a
la cible. Pour chaque alli\'e en position $\vec{g}$, on projette
$\vec{r}=\vec{g}-\vec{m}$ (depuis le canon $\vec{m}$) :
\[
\text{proj} = \vec{r}\cdot\vec{d},\qquad
\text{couloir bloqu\'e si } 0 < \text{proj} < D \ \text{et}\
\big\lVert \vec{r} - \text{proj}\,\vec{d}\big\rVert < \texttt{ALLY\_LANE\_WIDTH}=26.
\]
Si le couloir est bloqu\'e, le garde ne tire pas et se repositionne.

\subsection{4.8 Coordination de Groupe}
\begin{itemize}[leftmargin=*]
  \item \textbf{Alerte locale} : en entrant en combat, un garde pr\'evient ceux situ\'es
        dans \param{ALERT\_RADIUS}{420}\,px (pas tout le b\^atiment).
  \item \textbf{Le bruit des tirs alerte} : chaque coup de feu r\'eveille les gardes
        \`a moins de \param{GUNFIRE\_ALERT\_RADIUS}{520}\,px, qui viennent investiguer.
  \item \textbf{M\'emoire partag\'ee} : la derni\`ere position connue est diffus\'ee via
        \code{GameState.report\_player\_pos}.
\end{itemize}

\subsection{4.9 Recherche Coordonn\'ee}
En \'etat SEARCH, chaque garde inspecte \emph{sa} zone autour d'une ancre propre,
plac\'ee sur une position \textbf{anticip\'ee de la fuite} :
\[
\vec{ancre} = \vec{p}_{\text{dernier vu}} + \underbrace{\vec{v}_{\text{joueur}}\cdot \texttt{SEARCH\_LEAD}}_{\text{anticipation}}
+ \texttt{SEARCH\_SPREAD}\,(\cos\alpha_i,\ \sin\alpha_i),
\]
o\`u $\alpha_i$ est un angle stable et unique par garde (dispersion). Cela couvre
le terrain au lieu de provoquer un attroupement. \param{SEARCH\_LEAD}{0.55}\,s,
\param{SEARCH\_SPREAD}{150}\,px, \param{SEARCH\_RADIUS}{120}\,px.

% ============================== 5. NIVEAUX ==============================
\section{Design des Niveaux}

\subsection{5.1 Syst\`eme Data-Driven}
Les niveaux sont d\'ecrits dans \code{levels.gd} (\code{LevelCatalog}) sous forme de
dictionnaires \{ \texttt{name, subtitle, theme, ambiance, needs\_objective, map,
guards, medkits, terminals} \}. Le g\'en\'erateur \code{level.gd} lit
\code{LevelCatalog.get\_level(GameState.current\_level)} et construit la sc\`ene
selon le th\`eme. Chaque carte est une grille ASCII $46\times32$, $\texttt{TILE}=64$\,px.

\begin{center}
\begin{tabular}{cl}
\toprule
\textbf{Symbole} & \textbf{Signification} \\
\midrule
\texttt{\#} & obstacle plein (mur / lisi\`ere d'arbres / immeuble) \\
\texttt{o} & couverture basse (caisse / rocher / voiture) \\
\texttt{W} & eau (infranchissable, bloque aussi la vue) \\
\texttt{.} & sol praticable \\
\texttt{P} & entr\'ee joueur \quad \texttt{C} coffre (objectif) \quad \texttt{X} sortie/exfil \\
\bottomrule
\end{tabular}
\end{center}

\subsection{5.2 Niveau 0 --- Centre de Recherche}
Int\'erieur sombre (ambiance gris-bleu). Th\`eme \code{research} : tuiles textur\'ees,
mobilier (bureaux, serveurs, sofas, lasers). 8 gardes, 2 trousses de soins, 1
terminal, 1 coffre central. \textbf{Seul niveau \`a objectif obligatoire.}

\subsection{5.3 Niveau 1 --- La For\^et}
Ext\'erieur vert. Th\`eme \code{forest} : arbres et buissons (sprites + feuillages
proc\'eduraux), rochers. Une \textbf{rivi\`ere} (cases \texttt{W}, lignes 15--16)
coupe la carte avec un \textbf{pont} (passage praticable). 10 gardes,
2 trousses. Traverser du bas vers l'exfil en haut.

\subsection{5.4 Niveau 2 --- La Ville}
Th\`eme \code{city} : fond asphalte, immeubles proc\'eduraux, voitures, marquages
routiers. Un \textbf{fleuve} (lignes 3--4) avec un pont menant \`a l'\textbf{h\'elipad}
d'extraction (rep\`ere « H »). 10 gardes, 2 trousses.

\subsection{5.5 Navigation}
Pathfinding par \code{AStarGrid2D} : les cases \texttt{\#}, \texttt{o} et \texttt{W}
sont marqu\'ees solides ; d\'eplacement diagonal autoris\'e seulement sans obstacle.
Au chargement, \code{\_validate\_nav()} effectue un BFS depuis \texttt{P} et avertit
si le coffre, la sortie, une trousse, un terminal ou une case de garde est
inaccessible ou solide.

\subsection{5.6 R\`egles de Level Design}
\begin{itemize}[leftmargin=*]
  \item Au moins deux chemins entre zones (choix tactique).
  \item Couverture r\'eguli\`ere (jamais de longue ligne droite \`a d\'ecouvert).
  \item Le point d'exfiltration est lisible (rep\`ere visuel, h\'elipad en ville).
  \item \'Eclairage : zones sombres = s\^ures, zones \'eclair\'ees = expos\'ees.
  \item Connectivit\'e P$\to$X valid\'ee par BFS pour les trois cartes.
\end{itemize}

% ============================== 6. INTERFACE ==============================
\section{Interface et Audio}

\subsection{6.1 HUD}
CanvasLayer minimal : vies (c\oe urs), barre de sprint, munitions, libell\'e
d'alarme color\'e, score, minimap. La banni\`ere d'intro affiche
« NIVEAU $n$/3 --- \emph{nom} » avec sous-titre, puis s'estompe.

\subsection{6.2 \'Ecrans}
\begin{center}
\begin{tabular}{ll}
\toprule
\textbf{\'Ecran} & \textbf{Contenu} \\
\midrule
Menu principal & Logo, boutons Jouer / Quitter (Quitter masqu\'e sur mobile) \\
Briefing & Objectifs et consignes avant la mission \\
Jeu en cours & HUD minimal, immersion \\
Pause & Reprendre / Recommencer / Menu principal / Quitter \\
Mission r\'eussie & Score d\'etaill\'e (temps, gardes, discr\'etion, total) \\
Mission \'echou\'ee & Cause, Recommencer / Quitter \\
\bottomrule
\end{tabular}
\end{center}

\subsection{6.3 Sons}
Tir silencieux, tir de garde, rechargement, impact, garde abattu, clic UI,
exfiltration, plus une musique adaptative pilot\'ee par l'\'etat d'alarme
(\code{AudioManager}).

% ============================== 7. ARCHITECTURE ==============================
\section{Architecture Technique}

\subsection{7.1 Sc\`enes principales}
\begin{center}
\begin{tabular}{ll}
\toprule
\textbf{Sc\`ene} & \textbf{R\^ole} \\
\midrule
Main.tscn & Menu principal (point d'entr\'ee). \\
Briefing.tscn & \'Ecran de briefing $\to$ lance la campagne. \\
Level\_01.tscn & Sc\`ene g\'en\'erique recharg\'ee pour chaque niveau (data-driven). \\
Player.tscn & CharacterBody2D, Sprite, Muzzle, CoverArea. \\
Guard.tscn & CharacterBody2D, machine \`a \'etats, c\^one de vision. \\
Bullet.tscn & Area2D rapide, collisions par couche. \\
HUD.tscn, TouchControls.tscn & Interface et double-stick tactile. \\
Coffre / Terminal / Medkit / Exit & Objets interactifs (Area2D). \\
\bottomrule
\end{tabular}
\end{center}

\subsection{7.2 Autoloads}
\code{GameState} (\'etat d'alarme, score, vies, progression de campagne, signaux)
et \code{AudioManager} (musique adaptative, SFX), persistants \`a travers les
rechargements de sc\`ene.

\subsection{7.3 Couches de Collision}
1=\emph{world}, 2=\emph{player}, 3=\emph{guards}, 4=\emph{bullets\_player},
5=\emph{bullets\_enemy}, 6=\emph{interactables}. La couche 1 bloque \`a la fois le
mouvement \emph{et} la vision (= couverture).

\subsection{7.4 Assets et Licences}
Sprites Kenney.nl (CC0) et pack \code{mes\_assets} (nature, v\'ehicules, routes,
256\,px top-down). Aucune ressource sous licence restrictive.

% ============================== 8. DEPLOIEMENT ==============================
\section{D\'eploiement Android / Play Store}

\subsection{8.1 Configuration Mobile}
Orientation \emph{landscape} verrouill\'ee ; \code{stretch = canvas\_items} +
\code{aspect = expand} (adaptation \`a tous les ratios) ; renderer
\code{gl\_compatibility} (mobile) ; \'emulation tactile/souris crois\'ee ;
\code{quit\_on\_go\_back = false} pour g\'erer le bouton Retour (pause au lieu de
quitter).

\subsection{8.2 Presets d'Export}
Deux presets dans \code{export\_presets.cfg} : « Android APK » (test rapide,
\emph{runnable}) et « Android AAB (Play Store) » (publication). Communs :
package \code{com.eclipse.runaway}, version \code{1.0} (code 1), build Gradle,
architectures \code{arm64-v8a} + \code{armeabi-v7a}, mode immersif, \textbf{aucune
permission} (jeu hors-ligne).

\subsection{8.3 Signature}
Cl\'e release g\'en\'er\'ee hors d\'ep\^ot par \code{tools/gen\_keystore.ps1} et lue par
Godot via les variables d'environnement
\code{GODOT\_ANDROID\_KEYSTORE\_RELEASE\_*}. Le \code{.gitignore} exclut
\code{*.keystore}, \code{build/}, \code{*.aab}, \code{*.apk}.

\subsection{8.4 Conformit\'e Play Store}
AAB obligatoire (configur\'e), API cible 35, ic\^ones lanceur + adaptive + listing
512\,px g\'en\'er\'ees (\code{mobile\_icons/}). \`A fournir c\^ot\'e console : captures
d'\'ecran, feature graphic 1024$\times$500, politique de confidentialit\'e (aucune
donn\'ee collect\'ee), questionnaire de contenu (violence cartoon $\to$ PEGI~12 /
Teen). \textbf{Incr\'ementer \code{version/code} \`a chaque mise \`a jour.}

% ============================== 9. GLOSSAIRE ==============================
\section{Glossaire}
\begin{description}[leftmargin=2.5cm,style=nextline]
  \item[Top-Down] Vue de dessus ; la cam\'era regarde le personnage depuis le haut.
  \item[AAB] Android App Bundle ; format de publication exig\'e par le Play Store.
  \item[AStarGrid2D] Pathfinding sur grille utilis\'e par les gardes.
  \item[Suite de Weyl] $\operatorname{frac}(n\cdot\phi)$ ; r\'epartit des angles de
        fa\c{c}on quasi-uniforme (encerclement des gardes).
  \item[Lead targeting] Vis\'ee anticip\'ee tenant compte du temps de vol du projectile.
  \item[Hyst\'er\'esis] Zone morte \'evitant les oscillations de positionnement.
  \item[i-frame] Fen\^etre d'invincibilit\'e temporaire apr\`es un coup.
  \item[Data-driven] Niveaux d\'ecrits par des donn\'ees (catalogue) plut\^ot qu'en dur.
\end{description}

\vfill
\begin{center}\small\itshape
Document g\'en\'er\'e \`a partir du code source de production --- RUNAWAY : Operation Eclipse, v4.0.
\end{center}

\end{document}
